% Section 4.6.2: Qwen Within-Family Scale Sweep
% To be integrated into experiments.tex

\subsubsection{Within-Family Scale Sweep: Qwen 3B$\to$7B$\to$14B$\to$32B}
\label{sec:qwen_scale_sweep}

To test whether the suggestive scale trend observed under cross-family extraction (Section~\ref{sec:cross_family_equalized}) replicates within a single model family, we conducted equalized evaluations on four Qwen~2.5 models spanning 3B to 32B parameters ($n=100$ each). This provides the first comprehensive within-family scale analysis for behavioral deception detection, addressing the reviewer concern that prior evidence rested on a single non-significant within-family test (Qwen 7B$\to$14B, Fisher $p=0.17$).

\paragraph{Method.}
Following the prompt-equalized protocol (Section~\ref{sec:prompt_equalized}), we evaluated Qwen~2.5 models at four scales: 3B, 7B, 14B, and 32B via local Ollama with Claude Haiku~4.5 as interrogator/extractor (cross-family). All evaluations used identical neutral prompts, matched factual claims (true vs.\ modified-false), and adaptive interrogation with refusal-count LOO analysis.

\paragraph{Results.}
Table~\ref{tab:qwen_scale_sweep} shows refusal-count LOO accuracy across the four Qwen models. The pattern is \textbf{non-monotonic}: accuracy improves from 3B (50\%, chance) to 14B (72\%, peak), then collapses at 32B back to 50\% (chance).

\begin{table}[h]
\centering
\caption{Qwen within-family scale sweep. Refusal-count LOO accuracy with Fisher exact tests for adjacent increments. Mean refusal counts computed from interrogation transcripts using regex patterns (Section~\ref{sec:deployment}).}
\label{tab:qwen_scale_sweep}
\begin{tabular}{lcccccc}
\toprule
\textbf{Model} & \textbf{n} & \textbf{LOO} & \textbf{Ref (T)} & \textbf{Ref (L)} & \textbf{vs.\ Prev} & \textbf{Fisher $p$} \\
\midrule
Qwen 2.5 3B  & 100 & 50.0\% & 0.32 & 0.22 & -- & -- \\
Qwen 2.5 7B  & 100 & 68.0\% & 0.10 & 0.52 & +18\,pp & \textbf{0.014}* \\
Qwen 2.5 14B & 100 & 72.0\% & 0.18 & 0.76 & +4\,pp & 0.644 \\
Qwen 2.5 32B & 100 & 50.0\% & 0.00 & 0.00 & $-$22\,pp & \textbf{0.002}* \\
\bottomrule
\end{tabular}
\end{table}

The 3B$\to$7B increment is statistically significant (Fisher exact $p=0.014$), marking the \textbf{first significant within-family scale improvement} observed in behavioral deception detection. The 7B$\to$14B increment continues the trend (+4\,pp, $p=0.644$ non-significant), reaching peak accuracy at 14B. However, the 14B$\to$32B increment shows a significant \emph{decline} ($p=0.002$), returning to chance-level performance.

\paragraph{Behavioral shift at 32B.}
The 32B collapse is explained by a qualitative behavioral change: mean refusal count drops to \textbf{0.00 for both truth and lie conditions}. The model produces zero correction/refusal markers in interrogation transcripts, rendering the refusal-count heuristic ineffective. Manual inspection confirms that Qwen~32B knows the claims are false (a direct knowledge pre-check yields correct answers on 92\% of claims) but suppresses correction language entirely; responses are uniformly smooth and non-corrective regardless of claim truthfulness, indicating RLHF training optimized for agreeableness rather than knowledge-conflict signaling. This contrasts with 7B/14B models, which produce detectable correction language when confronted with false claims (mean refusal: 0.52--0.76 in lying trials).

\paragraph{Interpretation.}
The U-shaped pattern (50\%$\to$68\%$\to$72\%$\to$50\%) demonstrates that scale effects in behavioral deception detection are \textbf{non-monotonic} and \textbf{family-specific}. Three implications:

\textbf{First}, simple refusal-count heuristics achieve peak performance at \textbf{intermediate scale} (14B: 72\%), not at the largest scale tested (32B: 50\%). The assumption that ``larger is better'' does not hold when behavioral signals depend on RLHF training objectives that vary across scales.

\textbf{Second}, the significant 3B$\to$7B improvement ($p=0.014$) validates that scale \emph{can} improve equalized detection within a single family, addressing the concern that prior evidence rested on cross-family comparisons confounded by family composition (Section~\ref{sec:cross_family_equalized}). However, the 14B$\to$32B collapse demonstrates that this improvement is not monotonic and depends on model training rather than parameter count alone.

\textbf{Third}, the 32B behavioral shift (zero refusal markers) indicates a \textbf{limit of surface-level detection approaches}. When models are trained to avoid correction/refusal language---plausibly to optimize user satisfaction or helpfulness metrics---simple lexical heuristics fail regardless of model capability. This suggests that effective behavioral detection at the largest scales may require moving beyond refusal-marker counting to more sophisticated feature extraction or multi-dimensional scoring, though the current work does not test this hypothesis.

\paragraph{Comparison to cross-family trend.}
The cross-family trend reported in Section~\ref{sec:cross_family_equalized} shows monotonic improvement under Mistral Large extraction: Llama~3B (52\%) $\to$ Mistral~7B (62\%) $\to$ Llama~70B (69\%). This contrasts with the Qwen within-family U-shape. The discrepancy likely reflects:

\begin{enumerate}[leftmargin=*,itemsep=2pt]
\item \textbf{Family composition confounding}: Cross-family comparison mixes Llama and Mistral models with different training regimes. Llama~70B may produce more correction language than Qwen~32B due to different RLHF objectives.
\item \textbf{Generation gaps}: Llama~3.1/3.3 vs.\ Qwen~2.5 represents different model generations with evolving safety training.
\item \textbf{Extractor sensitivity}: Mistral Large may assign non-zero correction scores to Llama~70B transcripts where Qwen~32B would receive zero, masking the behavioral shift.
\end{enumerate}

The within-family analysis thus provides \textbf{cleaner evidence} by controlling for family-specific training differences, revealing that scale effects depend critically on RLHF objectives rather than parameter count alone.

\paragraph{Llama within-family replication (3B$\to$8B$\to$70B).}
To validate whether the Qwen U-shaped pattern generalizes to other model families, we replicated the within-family scale analysis on Llama models: 3.2~3B, 3.1~8B, and 3.3~70B, all evaluated under the prompt-equalized protocol with refusal-count LOO analysis ($n=100$ each, Claude Haiku~4.5 cross-family extraction).

Table~\ref{tab:llama_scale_sweep} shows the Llama results. Unlike Qwen's U-shape (50\%$\to$68\%$\to$72\%$\to$50\%), Llama exhibits a \textbf{flat-then-jump pattern}: no improvement from 3B to 8B (both 57.0\%, Fisher $p=1.000$), followed by a significant jump to 70B (76.0\%, $+$19\,pp, $p=0.004$).

\begin{table}[h]
\centering
\caption{Llama within-family scale sweep. Refusal-count LOO accuracy with Fisher exact tests for adjacent increments.}
\label{tab:llama_scale_sweep}
\begin{tabular}{lcccccc}
\toprule
\textbf{Model} & \textbf{n} & \textbf{LOO} & \textbf{Ref (T)} & \textbf{Ref (L)} & \textbf{vs.\ Prev} & \textbf{Fisher $p$} \\
\midrule
Llama 3.2 3B  & 100 & 57.0\% & 2.18 & 2.42 & -- & -- \\
Llama 3.1 8B  & 100 & 57.0\% & 1.56 & 2.56 & 0\,pp & 1.000 \\
Llama 3.3 70B & 100 & 76.0\% & 0.84 & 2.32 & +19\,pp & \textbf{0.004}* \\
\bottomrule
\end{tabular}
\end{table}

\paragraph{Family-specific scale effects.}
Comparing Qwen and Llama reveals family-specific patterns:

\textbf{Qwen (U-shape):} 50\% (3B) $\to$ 68\% (7B, $p=0.014$*) $\to$ 72\% (14B) $\to$ 50\% (32B, collapse). Peak at intermediate scale, driven by RLHF agreeableness training removing refusal markers at 32B.

\textbf{Llama (flat-then-jump):} 57\% (3B) $\to$ 57\% (8B, no change) $\to$ 76\% (70B, $p=0.004$*). Significant improvement only at largest scale, with no 32B collapse observed.

Three implications:

\textbf{First}, scale effects are \textbf{family-specific}, not universal. The Qwen U-shape does not replicate on Llama. Different model families exhibit different scale-dependent patterns, likely driven by family-specific RLHF training objectives and data.

\textbf{Second}, both families show at least one \textbf{significant within-family increment} (Qwen 3B$\to$7B: $p=0.014$; Llama 8B$\to$70B: $p=0.004$), validating that scale \emph{can} improve equalized detection within a single family. This addresses the concern that prior cross-family evidence was confounded by family composition.

\textbf{Third}, the absence of a 32B collapse in Llama (unlike Qwen) demonstrates that the zero-refusal-marker phenomenon is training-regime-specific, not an inevitable consequence of scale. Llama 70B (generation 3.3, released late 2024) retains correction behavior under knowledge conflict, while Qwen 32B (generation 2.5) was trained for agreeableness. Future models may exhibit similar collapse if RLHF objectives prioritize user satisfaction over knowledge-conflict signaling.

\paragraph{Limitations.}
This analysis tests two model families (Qwen at 4 scales, Llama at 3 scales). Generalization to other families (e.g., Mistral, Gemma, DeepSeek) remains unknown. The Llama analysis spans three model \emph{generations} (3.2, 3.1, 3.3), introducing potential generational confounds beyond parameter scale. Ideally, a within-generation sweep (e.g., Llama 3.3 at 1B/3B/8B/70B) would isolate scale from architectural changes, but such data is unavailable. Additionally, we test only refusal-count detection; more sophisticated behavioral features may not exhibit family-specific patterns, though EXP-J (Section~\ref{sec:deployment}) found refusal count matched or exceeded the full LLM pipeline on 6 of 7 models, suggesting limited upside from feature complexity.
